\section{Conclusion}

%Modeling an object is a tiresome and time expensive process by nature. %
%Creating variants is all the more so, considering that it requires starting over
%the process as many times as necessary. %
%Reevaluation mechanisms are, nowadays, increasingly frequent in parametric
%systems and offer an easier way to generate variants of an object. %
%A long-standing issue in CAD's field, called persistent naming, is to make sure
%the operations' parameters are still valid despite their modification. %
%To our knowledge, there does not seem to be such a mechanism taking advantage of
%graph transformation rules available so we propose one here. %

In this paper, we widen the naming problem studies to the rule-based graph transformation modeling systems. %
We take advantage of the formalism of both generalized maps and graph transformation rules  to tackle the reevaluation mechanism task. %
Ge\-ne\-ra\-li\-zed maps offer an homogeneous representation of an object in all
dimensions while Jerboa's rules %make possible the definition of 
define
geometric
modeling o\-pe\-ra\-tions on which it is actually possible to perform syntactical
a\-na\-ly\-sis. %
We implement: a persistent name scheme where each persistent name represents a unique 
dart's history through the successive %involved applications %
applications of rules 
and their matching nodes. 
%Thus, it enables to designate an orbit incident to a dart. %
Then, for each per\-sis\-tent name, an evaluation DAG
is built in order to trace an orbit's history from the bottom and up to the
orbits it originates from. %
To our knowledge, unlike other methods, our solution tracks only the entities used in the parametric specification and the ones they o\-ri\-gi\-na\-te from.
Representing the complete history of an orbit in an e\-va\-lua\-tion DAG allows for an
efficient persistent naming mechanism that takes into account the impact of both
origins and traces mo\-di\-fi\-ca\-tions during reevaluation. %
Finally, reevaluation DAGs are built from a top-down traversal of evaluation DAGs and
allow matching each topological parameter on one or more,  sometimes none, va\-lues depending on the editing of the parametric specification. %
%In our solution 
Thanks to our method,
not only the naming problem is tackled within the usual framework of parameters edition, but we also take the spe\-ci\-fi\-ca\-tion edition into account (operation addition, deletion or move).
Moreover, this approach provides all the possible updated values of the parameters and, thus, enables implementing different strategies.

More complex operations can make use of several rules 
brought together
in a script. %
Later works will revolve around widening this reevaluation mechanism
to scripts. %